% The 30 of 308 simulated excitation runs excluded by the onset-free
% run gates, grouped by cause.  Classification: full-window slope and
% linearity recomputed per run; "developed portion" = samples above
% the slope-gate moment floor.  Reproduced by the gate audit in the
% analysis notes (extract_piecewise_batch gate criteria).
\begin{table}[t]
\caption{Simulated excitation runs excluded by the run gates
($30$ of $308$; the identification uses the remaining $278$).}
\label{tab:sim_gate_exclusions}
\centering
\begin{tabular}{@{}p{0.26\linewidth}cp{0.42\linewidth}@{}}
\toprule
Group & Runs & Signature \\
\midrule
Actuator transient at fast ramps
($\dot M \ge 0.65$\,N$\cdot$m/s)
 & 20
 & Window shortens toward $\tau_{\rm eff} \approx 83$\,ms, so rotor
   spin-up curvature enters it: linearity $57$--$221$\,mN$\cdot$m;
   the developed portion still tracks the commanded rate within
   $3\%$.\\[2pt]
Slope-floor artifact on boundary cases
 & 8
 & Well-executed ramps (full-window slope within $1\%$, linearity
   $\approx 8$\,mN$\cdot$m) whose window ends barely above the
   slope-gate moment floor, leaving the masked estimator too few
   samples --- conservative false rejections.\\[2pt]
Threshold-marginal
 & 2
 & Ramp-rate error $3.2\%$ against the $3\%$ gate (S9/$M_x$ at
   $0.20$ and $0.30$\,N$\cdot$m/s).\\
\bottomrule
\end{tabular}\\[2pt]
\raggedright\footnotesize The residual bound of
\cref{subsubsec:sim_rmse_bound} holds on the rejected runs as well:
the gates protect the accuracy of the critical-moment readout, not
the validity of the bound.
\end{table}
